\section*{Repository Appendix B: Detailed audit methods and secondary diagnostics}

\label{supp:audit_details}

\subsection*{B.1 Detailed confirmation and abstention audit design}

Fixed-weight confirmation was examined through four complementary audits addressing robustness, bootstrap validity, multiplicity, and possible optimism from candidate selection. These audits did not modify the benchmark gate or the evidence taxonomy. They provided additional checks on candidates that had already completed fixed-weight validation.

The first audit examined the robustness properties implied by the frozen weight vectors. For each confirmed specialist, it recorded the weights assigned to the sample mean, median, and modal components and interpreted them through the influence-function decomposition introduced in Section~2 of the main text. A positive weight on the sample mean implies a formal breakdown point of zero, even when most remaining weight is assigned to estimators with bounded influence. Gross-error sensitivity was assessed relative to the sample mean by examining how the linear term associated with $w_{\mathrm{mean}}$ scales the effect of an arbitrarily large observation. Geometric and harmonic components were considered separately because their contributions grow more slowly than the linear contribution of the sample mean.

The second audit addressed bootstrap validity for nonsmooth components. The ordinary bootstrap is unreliable for the mode and may converge slowly for the median~\citep{ghosh1984median,romano1988mode}. The practical relevance of this limitation was assessed by inspecting the total weight assigned to modal and median components in each confirmed or otherwise important candidate. The modal weights of the confirmed specialists were very small, whereas candidates with more substantial median weight were retained with an explicit interpretive caveat. An $m$ out of $n$ subsampling interval was identified as an appropriate extension for future confirmation studies.

The third audit examined multiplicity across confirmation candidates. The Benjamini--Hochberg procedure was applied to the candidate-level confirmation signals at a 5\% false discovery rate. This correction was an additional audit and did not replace the original dual benchmark gate.

The fourth audit evaluated possible optimism from selecting candidates before reporting their gains. The eight validation seeds were divided into two groups: one half was used for selection and the other for performance reporting. These split-sample estimates were compared with the pooled estimates from all eight seeds. Because each half contained only four seeds, the results were treated as robustness estimates, not as separate inferential intervals.

A complementary audit examined benchmark-retained outcomes. Its purpose was to distinguish genuine abstention from a failure of the genetic algorithm to locate an existing gate-passing region. All benchmark-retained regimes from the expanded discovery stage were evaluated through an independent random Dirichlet search over the same 26-component simplex. The audit covered 17 regimes and both validation modes, producing 34 regime and mode cells.

For each regime and validation mode, the audit generated 4,000 random weight vectors from a symmetric Dirichlet distribution with concentration parameter $\alpha=0.30$. Every draw was evaluated with the same simulator, admissibility rules, benchmark registry, definition of $q_{0.95}$, and Monte Carlo budget used in fixed-weight validation, with $R=500$. The original dual gate required improvement in both MSE and $q_{0.95}$, together with stability across all eight validation seeds,
\[
303,\ 404,\ 505,\ 606,\ 707,\ 808,\ 909,\ 1010.
\]

A regime and mode cell confirmed abstention when no random draw cleared the dual gate at all eight seeds. The frozen weights of the confirmed transfer specialists were included as positive controls to verify that the audit could recover weight vectors already known to satisfy the gate. Randomly identified passing vectors were not treated as confirmed discoveries. They were recorded only as targets for a future search in which their weights would need to be persisted, frozen, and subjected to the complete validation procedure.

\subsection*{B.2 Seed-specific fixed-weight confirmation results}

\begin{table}[htbp]
\centering
\caption{Seed-specific fixed-weight confirmation results for the three
Lognormal candidates. Mean and $q_{0.95}$ gains are relative to the
criterion-specific strongest admissible benchmarks within each validation
seed. The $q_{0.95}$ interval is the 95\% paired bootstrap interval based on
$B=500$ resamples for that seed. A dual-gate pass requires positive gains
in both MSE and $q_{0.95}$.}
\label{tab:seed_specific_confirmation}
\scriptsize
\setlength{\tabcolsep}{2.3pt}
\begin{tabular}{@{}P{0.12\textwidth}P{0.12\textwidth}P{0.07\textwidth}
P{0.12\textwidth}P{0.12\textwidth}P{0.20\textwidth}P{0.09\textwidth}@{}}
\toprule
\textbf{Candidate} &
\textbf{Mode} &
\textbf{Seed} &
\textbf{Mean gain} &
\textbf{$q_{0.95}$ gain} &
\textbf{$q_{0.95}$ paired bootstrap 95\% CI} &
\textbf{Dual gate} \\
\midrule
CV-019 & Locked unseen & 303 & 5.18\% & 18.36\% & [16.36\%, 20.12\%] & Pass \\
CV-019 & Locked unseen & 404 & 6.15\% & 17.35\% & [12.35\%, 23.20\%] & Pass \\
CV-019 & Locked unseen & 505 & 7.90\% & 13.16\% & [7.62\%, 18.55\%] & Pass \\
CV-019 & Locked unseen & 606 & 4.20\% & 13.73\% & [10.40\%, 16.79\%] & Pass \\
CV-019 & Locked unseen & 707 & 6.88\% & 18.64\% & [16.11\%, 19.72\%] & Pass \\
CV-019 & Locked unseen & 808 & 6.13\% & 18.48\% & [16.63\%, 19.74\%] & Pass \\
CV-019 & Locked unseen & 909 & 7.29\% & 16.78\% & [8.37\%, 20.98\%] & Pass \\
CV-019 & Locked unseen & 1010 & 6.37\% & 13.96\% & [8.97\%, 18.79\%] & Pass \\
\addlinespace[0.35em]
CV-010 & Locked unseen & 303 & 1.93\% & 2.21\% & [1.95\%, 2.72\%] & Pass \\
CV-010 & Locked unseen & 404 & 2.08\% & 2.02\% & [1.73\%, 2.44\%] & Pass \\
CV-010 & Locked unseen & 505 & 1.63\% & 2.13\% & [1.57\%, 2.65\%] & Pass \\
CV-010 & Locked unseen & 606 & 2.23\% & 2.37\% & [1.95\%, 2.73\%] & Pass \\
CV-010 & Locked unseen & 707 & 2.03\% & 2.11\% & [1.80\%, 2.68\%] & Pass \\
CV-010 & Locked unseen & 808 & 2.17\% & 2.30\% & [1.96\%, 2.80\%] & Pass \\
CV-010 & Locked unseen & 909 & 1.67\% & 2.53\% & [2.05\%, 3.11\%] & Pass \\
CV-010 & Locked unseen & 1010 & 1.99\% & 2.27\% & [1.78\%, 2.62\%] & Pass \\
\addlinespace[0.35em]
CV-011 & Original & 303 & 0.52\% & 3.41\% & [2.75\%, 4.05\%] & Pass \\
CV-011 & Original & 404 & -0.45\% & 3.67\% & [2.89\%, 4.37\%] & Retain \\
CV-011 & Original & 505 & 0.92\% & 4.45\% & [3.62\%, 5.04\%] & Pass \\
CV-011 & Original & 606 & -0.06\% & 3.85\% & [3.26\%, 4.38\%] & Retain \\
CV-011 & Original & 707 & -0.25\% & 3.57\% & [2.90\%, 3.92\%] & Retain \\
CV-011 & Original & 808 & -0.20\% & 3.91\% & [3.34\%, 4.62\%] & Retain \\
CV-011 & Original & 909 & -0.65\% & 3.72\% & [2.98\%, 4.40\%] & Retain \\
CV-011 & Original & 1010 & -0.55\% & 3.98\% & [3.42\%, 4.55\%] & Retain \\
\bottomrule
\end{tabular}
\end{table}
\FloatBarrier

\subsection*{B.3 Seed-level confirmation matrix}

\begin{figure}[htbp]
\centering
\includegraphics[width=.90\textwidth]{fig_confirmatory_seed_survival.png}
\caption{Seed-level expanded-gate survival for the confirmed and local Lognormal signals. Each cell indicates whether a candidate satisfied the expanded benchmark gate for the corresponding validation seed.}
\label{fig:seed_survival}
\end{figure}

\subsection*{B.4 Expanded-basis evidence taxonomy by family}

\begin{figure}[htbp]
\centering
\includegraphics[width=0.95\textwidth]{fig_phase2_taxonomy_publication.png}
\caption{Evidence taxonomy from post-discovery fixed-weight validation ($N_{\mathrm{EST}}=26$), shown by family. Transfer specialists have interval-supported advantages in locked-unseen evaluation. Discovery-supported local candidates retain valid discovery evidence, near-gate candidates identify borderline regions for diagnosis, and negative-control successes verify appropriate benchmark retention.}
\label{fig:phase2_taxonomy_publication}
\end{figure}

\subsection*{B.5 Candidate-level taxonomy and gain matrix}

\begin{sidewaystable}[p]
\centering
\caption{Candidate-level evidence and gain matrix for the complete
25-candidate expanded-basis validation panel. Positive values indicate
improvement and negative values indicate loss relative to the
criterion-specific strongest admissible benchmark. Each gain cell reports
MSE gain followed by $q_{0.95}$ gain. Discovery values are the gains that
defined candidate status during expanded search. Locked-unseen and original
values are frozen-weight averages across the eight validation seeds. The CI
columns report the number of seeds whose $q_{0.95}$ paired bootstrap interval was
entirely positive. ``Transfer'' denotes transfer specialist;
``Discovery-local'' denotes an accepted discovery winner not confirmed by
the frozen gate; the remaining grades identify diagnostic near-gate and
negative-control cells.}
\label{tab:taxonomy_gain_matrix}

\tiny
\setlength{\tabcolsep}{1.6pt}

\resizebox{\textheight}{!}{%
\begin{tabular}{@{}lllllllll@{}}
\toprule
\textbf{Candidate} &
\textbf{Family / regime / seed} &
\textbf{Source} &
\textbf{Taxonomy} &
\textbf{Discovery mean / $q_{0.95}$} &
\textbf{Locked mean / $q_{0.95}$} &
\textbf{Locked $q_{0.95}$ CI seeds} &
\textbf{Original mean / $q_{0.95}$} &
\textbf{Original $q_{0.95}$ CI seeds} \\
\midrule
Q1R011 / FWVR011 & Weibull / REG01 / 101 & GA win & Transfer & +8.30\% / +7.32\% & +2.79\% / +2.72\% & 8/8 & -15.77\% / -45.06\% & 0/8 \\
Q1R012 / FWVR012 & Weibull / REG02 / 202 & GA win & Transfer & +29.06\% / +23.61\% & +3.18\% / +1.25\% & 8/8 & -15.20\% / -34.56\% & 0/8 \\
\addlinespace[0.35em]
FWVR001 & Exwald / REG01 / 101 & GA win & Discovery-local & +4.69\% / +10.44\% & -11.06\% / -7.36\% & 0/8 & -21.89\% / -48.44\% & 0/8 \\
FWVR002 & Exwald / REG03 / 101 & GA win & Discovery-local & +23.73\% / +11.49\% & -32.20\% / -58.58\% & 0/8 & -30.24\% / -57.94\% & 0/8 \\
FWVR003 & Exwald / REG03 / 202 & GA win & Discovery-local & +24.78\% / +20.62\% & -3.29\% / -8.37\% & 0/8 & -18.03\% / -40.42\% & 0/8 \\
FWVR004 & Invgauss / REG01 / 101 & GA win & Discovery-local & +52.21\% / +38.20\% & -8.43\% / +1.46\% & 0/8 & -36.22\% / -90.21\% & 0/8 \\
FWVR005 & Invgauss / REG02 / 101 & GA win & Discovery-local & +32.15\% / +16.94\% & -3.18\% / +2.62\% & 0/8 & -57.22\% / -164.93\% & 0/8 \\
FWVR006 & Invgauss / REG02 / 202 & GA win & Discovery-local & +17.53\% / +5.32\% & -5.52\% / +0.47\% & 0/8 & -42.91\% / -103.40\% & 0/8 \\
FWVR007 & Invgauss / REG03 / 202 & GA win & Discovery-local & +17.59\% / +8.34\% & -11.52\% / -1.58\% & 0/8 & -24.51\% / -50.99\% & 0/8 \\
FWVR008 & Lognormal / REG01 / 101 & GA win & Discovery-local & +16.76\% / +3.86\% & -38.07\% / -127.60\% & 0/8 & -24.76\% / -85.25\% & 0/8 \\
FWVR009 & Lognormal / REG03 / 101 & GA win & Discovery-local & +13.65\% / +3.82\% & -1.91\% / -0.72\% & 0/8 & -43.30\% / -164.67\% & 0/8 \\
FWVR010 & Normal / REG02 / 202 & GA win & Discovery-local & +1.68\% / +6.99\% & -0.09\% / -0.28\% & 0/8 & -0.03\% / -0.54\% & 0/8 \\
\addlinespace[0.35em]
FWVR013 & Exgaussian / REG01 / 101 & Near-gate & Near-gate & -16.20\% / -4.87\% & -1.20\% / -1.26\% & 0/8 & -1.02\% / -1.06\% & 0/8 \\
FWVR014 & Exgaussian / REG02 / 101 & Near-gate & Near-gate & -9.84\% / -13.01\% & -0.95\% / -0.98\% & 0/8 & -0.88\% / -1.07\% & 0/8 \\
FWVR015 & Exgaussian / REG03 / 101 & Near-gate & Near-gate & -10.02\% / -8.77\% & -1.16\% / -1.89\% & 0/8 & -0.84\% / -1.10\% & 0/8 \\
FWVR016 & Normal / REG01 / 101 & Near-gate & Near-gate & -4.52\% / -4.05\% & -0.23\% / -0.30\% & 0/8 & -0.16\% / -0.30\% & 0/8 \\
FWVR017 & Normal / REG02 / 101 & Near-gate & Near-gate & -1.57\% / -4.51\% & -0.15\% / -0.42\% & 0/8 & -0.20\% / -0.51\% & 0/8 \\
FWVR018 & Weibull / REG02 / 101 & Near-gate & Near-gate & -0.64\% / -0.75\% & -75.95\% / -283.58\% & 0/8 & -52.45\% / -197.54\% & 0/8 \\
FWVR019 & Weibull / REG03 / 101 & Near-gate & Near-gate & -71.64\% / -9.71\% & +0.66\% / -0.15\% & 2/8 & -31.66\% / -106.49\% & 0/8 \\
\addlinespace[0.35em]
FWVR020 & Exgaussian / REG01 / 202 & Control & Negative control & -0.39\% / +0.40\% & -2.26\% / -4.56\% & 0/8 & -2.17\% / -4.55\% & 0/8 \\
FWVR021 & Exgaussian / REG02 / 202 & Control & Negative control & -1.25\% / -2.10\% & -0.68\% / -0.50\% & 0/8 & -0.95\% / -1.05\% & 0/8 \\
FWVR022 & Exgaussian / REG03 / 202 & Control & Negative control & -8.44\% / -0.66\% & -0.44\% / -0.60\% & 0/8 & -0.69\% / -0.71\% & 0/8 \\
FWVR023 & Normal / REG01 / 202 & Control & Negative control & +0.69\% / -0.50\% & -0.07\% / -0.16\% & 0/8 & -0.01\% / +0.07\% & 0/8 \\
FWVR024 & Normal / REG03 / 101 & Control & Negative control & -1.58\% / -2.46\% & -0.19\% / -0.50\% & 0/8 & -0.19\% / -0.24\% & 0/8 \\
FWVR025 & Normal / REG03 / 202 & Control & Negative control & -0.13\% / -0.21\% & -0.12\% / -0.25\% & 0/8 & -0.04\% / -0.01\% & 0/8 \\
\bottomrule
\end{tabular}%
}
\end{sidewaystable}
\FloatBarrier

\subsection*{B.6 Frozen-weight robustness audit}

\begin{table}[htbp]
\centering
\small
\caption{Frozen-weight robustness audit for the two confirmed Weibull transfer specialists. Both estimators substantially attenuate gross-error sensitivity relative to the sample mean. Their small positive sample mean weights preserve a formal breakdown point of zero.}
\label{tab:robustness_pub}
\begin{tabular}{@{}P{0.23\textwidth}P{0.11\textwidth}P{0.12\textwidth}P{0.11\textwidth}P{0.31\textwidth}@{}}
\toprule
\textbf{Specialist} & $w_{\text{mean}}$ & $w_{\text{median}}$ & $w_{\text{mode}}$ & \textbf{Breakdown / gross-error influence} \\
\midrule
Weibull REG01 (seed 101) & 0.016 & 0.179 & 0.007 & $0$ / $\approx 1/62$ of $\bar X_n$ \\
Weibull REG02 (seed 202) & 0.022 & 0.018 & 0.004 & $0$ / $\approx 1/46$ of $\bar X_n$ \\
\bottomrule
\end{tabular}
\end{table}

\subsection*{B.7 Random-Dirichlet abstention audit}

\begin{table}[htbp]
\centering
\small
\caption{Results of the random Dirichlet abstention audit over the 17 benchmark-retained regimes. The audit comprised 4,000 draws per regime and validation mode, with $\alpha=0.30$ and stability across eight seeds. The confirmed transfer specialists served as positive controls.}
\label{tab:abstain_audit}
\begin{tabular}{lc}
\toprule
\textbf{Quantity} & \textbf{Value} \\
\midrule
Abstained regimes audited & $17$ ($34$ regime and mode cells) \\
Regime and mode cells confirming abstention & $23 / 34$ \\
Unique regimes fully confirmed in both modes & $9 / 17$ \\
Unique regimes with a Dirichlet signal & $8 / 17$ \\
Normal / Ex Gaussian controls confirmed & $6 / 6$ \\
Positive control, transfer specialists & $8 / 8$ and $8 / 8$ \\
Strongest flagged regime, Q1R003 Ex Wald & $1.8\%$ pass rate \\
\bottomrule
\end{tabular}
\end{table}

\FloatBarrier
\subsection*{B.8 External empirical gate funnel}

\begin{figure}[H]
\centering
\includegraphics[width=0.78\textwidth]{fig_realworld_evidence_funnel_corrected.png}
\caption{Conjunctive stringency of the external primary gate. The figure annotation reporting 2,211 strong-win cells refers to condition-level rows with positive paired bootstrap intervals among 5,900 audit-ready gate rows. These rows include repeated conditions within the same dataset, specialist, and mode comparison and are not bootstrap replicates or independent evidence. They were aggregated into 809 profile-eligible comparisons; 256 contained at least one strong win supported by a confidence interval, 255 survived Benjamini--Hochberg correction at 5\%, and 257 survived at the 10\% sensitivity threshold. The 10\% set includes one additional strict win and one sensitivity-only row.}
\label{fig:rw_funnel_pub}
\end{figure}
\FloatBarrier
